\documentclass[a4paper]{amsart}

\usepackage{color}
\usepackage[colorlinks=true, citecolor=red]{hyperref}  %needs to come before amsrefs package to work with \cite{}*{}
\usepackage{amssymb,latexsym}
\usepackage[alphabetic]{amsrefs}
\usepackage[mathscr]{eucal}
\usepackage{pdfsync}
\usepackage{graphicx}
\usepackage{epstopdf}
\DeclareGraphicsRule{.tif}{png}{.png}{`convert #1 `dirname #1`/`basename #1 .tif`.png}
\usepackage[all]{xy}
\usepackage{titletoc}
%\usepackage{showkeys}  %options notref  notcite 

\input{rgb.tex}

\CompileMatrices

\theoremstyle{plain}
\newtheorem{theorem}{Theorem}[section]
\newtheorem{lemma}[theorem]{Lemma}
\newtheorem{proposition}[theorem]{Proposition}
\newtheorem{corollary}[theorem]{Corollary}

\theoremstyle{definition}
\newtheorem{definition}[theorem]{Definition}%[section]


\theoremstyle{remark}
\newtheorem{defn}[theorem]{Definition} 
\newtheorem{example}[theorem]{Example}
\newtheorem{remark}[theorem]{Remark}
\newtheorem{discussion}[theorem]{Discussion}

%\numberwithin{equation}{section}

%\newcommand{\abs}[1]{\lvert#1\rvert}

\DeclareMathOperator{\Lip}{Lip}
\DeclareMathOperator{\dist}{dist}
\DeclareMathOperator{\expected}{\mathbb{E}}
\DeclareMathOperator{\prb}{Prob}
\DeclareMathOperator{\spt}{spt}
\DeclareMathOperator{\lap}{\Delta}
\DeclareMathOperator{\Tr}{Tr}

\raggedbottom

\addtolength{\textwidth}{2cm}
\addtolength{\oddsidemargin}{-1cm}
\addtolength{\evensidemargin}{-1cm} 
\addtolength{\topmargin}{-1cm} 
\addtolength{\textheight}{2cm}

 
%%% BEGIN DOCUMENT
\begin{document}

\title{PDE's}  
\author{John E.\ Hutchinson}
\date{\today}

%\begin{abstract} 
%\end{abstract}

\maketitle

\tableofcontents


%%%%%%%%%%%%%%%%%%%%%%%%%%%%%%%%
%%%%%%%%%%%%%%%%%%%%%%%%%%%%%%%%
%%%%%%%%%%%%%%%%%%%%%%%%%%%%%%%%
%%%%%%%%%%%%%%%%%%%%%%%%%%%%%%%%
\section{\textbf{22/4/09: Dirichlet Problem}} 
Comments  about solving the Dirichlet energy mininisation problem using Sobolev spaces.  

In particular,   note that \emph{the full minimising sequence  converges (not just a subsequence), and the convergence is in the (strong) $ W^{1,2} $ sense}.

%%%%%%%%%%%%%%%%%%%%%%%%%%%%%%%
\subsection{Sobolev Spaces}
Suppose $ \Omega \subset \mathbb{R}^n $ where $ \Omega $ is a bounded domain (i.e.\ open and connected).


\begin{definition}
\begin{align*} 
W^{1,2}(\Omega) &= \{ u \in L^2(\Omega) : Du \in L^2(\Omega), \text{ where $ Du $   is the weak derivative}\},\\
W_0^{1,2}(\Omega) &= \emph{closure of $ C_0^1(\Omega) $ in $ W^{1,2}(\Omega)$} .
\end{align*} 
\end{definition} 
See \cite{GT}*{\S\S 7.3, 7.5}

\medskip
By ``$  Du $ is the weak derivative'' we mean for $ i=1,\dots n $ there exists $ v_i \in L^2(\Omega) $ s.t.\
\begin{equation} \label{wd}
\int_\Omega u\, D_i\phi = - \int_\Omega v_i \phi \quad  \forall \phi \in C_0^1(\Omega),
\end{equation}
in which case $ v_i $ is unique and is denoted by $D_i u$.  

 \begin{remark}[An expression for the $ W^{1,2}(\Omega) $ norm]
From~\eqref{wd}, if $ u \in L^2(\Omega) $ then $ D_i u $ exists and is in $ L^2(\Omega) $, i.e.\  $ u \in W^{1,2} (\Omega)$,  iff the right side of ~\eqref{wde}  is finite, in which case
\begin{equation} \label{wde}
\int_\Omega |D_iu|^2 = \sup \left\{ \int_\Omega u\, D_i\phi : \phi \in C_0^1(\Omega), \ \|\phi\|_{L^2} \leq 1 \right\}.
\end{equation}
\end{remark} 

\begin{remark}[$ W^{1,2}$ Inner  Product] Note that $ W_0^{1,2} (\Omega)$ is an inner product space with 
\[
(f,g)_{W^{1,2}} = \int_\Omega Df\cdot Dg .
\]
Similarly for $ W^{1,2}(\Omega) $ if we take equivalence classes modulo  constant functions.
 \end{remark} 
 
 \begin{remark}[Weak Convergence]   \label{rwc}  One says  \emph{$ u_n  \rightharpoonup u $ (weakly) in $ L^2 (\Omega)$} if $ \int_\Omega u_n \phi \to  \int_\Omega u \phi $ for all $ \phi $ in some dense subset of $ L^2(\Omega ) $. It follows
 \[
 \int_\Omega |u|^2 \leq \liminf  \int_\Omega |u_n|^2.
 \]
 
 One says \emph{$ u_n  \rightharpoonup u $ (weakly) in $ W^{1,2}(\Omega) $} if $ u_n  \rightharpoonup u $   in $ L^2 $ and $ D_{u_n}  \rightharpoonup D_u $   in $ L^2 (\Omega)$ \ for all $ i $. 
 
 It follows from \eqref{wd} and \eqref{wde} that if $  u_n \in W^{1,2}(\Omega) $, $ u_n \rightharpoonup u $ in $ L^2(\Omega) $ and $ \int_\Omega|Du_n|^2 $ is uniformly bounded by $ K $ say, then $ u \in W^{1,2}(\Omega)$, $ \int_\Omega |Du|^2 \leq K $, and $ u_n \rightharpoonup u $ in $ W^{1,2}(\Omega) $. 
  \end{remark} 
  
\begin{remark}[Membership in $ W_0^{1,2}(\Omega) $] \label{rmw} One can check that if $ u \in W^{1,2}(\Omega) $ then $ u \in W_0^{1,2}(\Omega) $ iff 
$\int_\Omega D_i u\, \phi = -\int_{\Omega} u\, D_i   \phi   $ for all $ \phi $ in some dense subset of $ W^{1,2} (\Omega)$.   It follows that membership in $ W_0^{1,2}(\Omega) $ is preserved under  $ W ^{1,2}(\Omega) $ weak convergence.
 \end{remark} 
  
  
  
  
%%%%%%%%%%%%%%%%%%%%%%%%%%%%%%%
\subsection{Dirichlet PDE Problem}  The classical Dirichlet PDE problem
\[
\lap u =  g \text{  in }\Omega, \quad u = 0 \text{ on } \partial \Omega,
\]
can be reformulated as follows:

\medskip
\emph{Suppose $ g \in L^2(\Omega) $.  Solve}
\begin{equation}\label{ddp}
\begin{gathered}
u \in W_0^{1,2} (\Omega), \\
\int_\Omega Du\,  D\phi = - \int_\Omega g\phi \quad \forall \phi \in W_0^{1,2} (\Omega).
\end{gathered} 
\end{equation} 

Using the inner product on $ W_0^{1,2} (\Omega)$ and the Riesz representation theorem there exists a unique $ u \in W_0^{1,2} (\Omega)
$ solving~\eqref{ddp}.  Moreover $  u \in W^{2,2} (\Omega)$, and higher regularity follows from higher regularity of $ g $, see \cite{GT}*{Ch 8}.

\bigskip 
 The more general classical problem 
\[
\lap u =  g \text{  in }\Omega, \quad u = f \text{ on } \partial \Omega,
\]
is reformulated as follows after setting $ w = u-f $:

\medskip 
\emph{Suppose $ g \in L^2(\Omega) $, $ f \in W^{2,2}(\Omega) $.  Solve}
\begin{equation}\label{dddp}
\begin{gathered}
w \in W_0^{1,2} (\Omega), \\
\int_\Omega Dw\,  D\phi = - \int_\Omega (g - \lap f) \phi \quad \forall \phi \in W_0^{1,2} (\Omega).
\end{gathered} 
\end{equation} 
This is then covered by the previous problem.  

One only needs $ f \in  W^{1,2}(\Omega) $, if one replaces $ \int_\Omega \lap f \, \phi $ by $ -\int_\Omega Df\, D\phi (\Omega)$.  And similarly,  one can also take $ g \in W^{-1,2}(\Omega)$, the dual of $  W^{1,2}(\Omega)$.





%%%%%%%%%%%%%%%%%%%%%%%%%%%%%%%
\subsection{Dirichlet Energy Minimisation Problem} We   solve the energy minimisation problem, examine  the minimising sequence, and relate the minimiser to the previous PDE problem.  

The classical energy minimisation problem is reformulated as follows.

\medskip 
\emph{Suppose $ f \in W^{1,2}(\Omega) $.  Show there exists $ u \in W^{1,2}(\Omega) $ which minimises $ \int_\Omega |Du|^2 $, subject to the constraint $ u-f \in W_0^{1,2}(\Omega) $.}

\medskip (The constraint is a formulation of the condition $ u=f $ on $ \partial \Omega $.)

\begin{proof}[Proof and Discussion] \mbox{}

\textbf{1}.  \emph{Minimising Sequence}:  Let $ \gamma $ be the inf of $ \int_\Omega |Du|^2 $ such that $ u-f \in W_0^{1,2}(\Omega) $.  Let $ (u_n )\subset  W^{1,2}(\Omega)$ be a minimising sequence.

\medskip \textbf{2}.  \emph{$ L^2 $ convergent subsequence}:  Since $ W^{1,2}(\Omega) \subset L^2(\Omega ) $ is a compact imbedding, there is a subsequence $ (u_n) $ such that $ u_n \to u $ (say) in $ L^2 (\Omega)$ and $ \int_\Omega |Du|^2 \downarrow \gamma $.

\medskip \textbf{3}.  \emph{$ W^{1,2} $ weakly convergent subsequence}:  
From Remark \ref{rwc}, $ u \in W^{1,2}(\Omega)$, $ \int_\Omega |Du|^2 \leq \gamma $, and $ u_n \rightharpoonup u $ in $ W^{1,2} (\Omega)$. 

\medskip \textbf{4}. \emph{$ u $ is a minimiser}:  From Remark~\ref{rmw}, $ u-f \in W_0^{1,2}(\Omega) $ and so $ \int_\Omega |Du|^2 = \gamma $
since we already know $ \int_\Omega |Du|^2 \leq \gamma $ from \textbf{3}.

\medskip \textbf{5}. \emph{Strong $ W^{1,2} $ convergence of the \underline{full} sequence}:  We   get this directly from \textbf{1}   by noting
\[
\int_\Omega |Du_n|^2 + \int_\Omega |Du_m|^2 = \frac{1}{2} \int_\Omega |Du_n - Du_m|^2 + 2 \int_\Omega \left|D \frac{u_n+u_m}{2}\right|^2.
\]
The left side terms both decrease to $ \gamma $,    the second right side term is $ \geq \gamma $, and so $ \int_\Omega |Du_n - Du_m|^2 \to 0  $. 
By Poincare, $ \int_\Omega | u_n -  u_m|^2 \to 0 $.  

Then $ u-f \in W_0^{1,2}(\Omega) $ since $ W_0^{1,2}(\Omega) $ is certainly closed under strong $  W ^{1,2}$ convergence.

Moreover, $ \int_\Omega |Du|^2 \leq \gamma $ by Fatou's lemma, and so $ \int_\Omega |Du|^2= \gamma $.

\medskip \textbf{6}. \emph{$ u $ satisfies Laplace's equation}:     Suppose $ \phi \in W_0^{1,2}(\Omega)$.  Because
\[
\int_\Omega |D(u+\phi)|^2 = \int_\Omega |Du|^2 + 2 \int_\Omega Du\,  D\phi + \int_\Omega |D\phi|^2,
\]
it follows $\int_\Omega Du\, D\phi = 0 $.  Since otherwise, as small linear terms dominate  quadratic terms, we could find $ \phi $ such that $ u+\phi $ has less energy than $ u $. 
\end{proof} 




%%%%%%%%%%%%%%%%%%%%%%%%%%%%%%%%
%%%%%%%%%%%%%%%%%%%%%%%%%%%%%%%%
%%%%%%%%%%%%%%%%%%%%%%%%%%%%%%%%
%%%%%%%%%%%%%%%%%%%%%%%%%%%%%%%%
\begin{thebibliography}{9} 

\bib{GT}{book}{
   author={Gilbarg, David},
   author={Trudinger, Neil S.},
   title={Elliptic partial differential equations of second order},
   series={Grundlehren der Mathematischen Wissenschaften [Fundamental
   Principles of Mathematical Sciences]},
   volume={224},
   edition={2},
   publisher={Springer-Verlag},
%   place={Berlin},
   date={1983},
   pages={xiii+513},
%   isbn={3-540-13025-X},
%   review={\MR{737190 (86c:35035)}},
}

\end{thebibliography}

\end{document} 

